\documentclass[11pt]{article}

\usepackage[margin=1in]{geometry}
\usepackage{amsmath,amssymb,booktabs,array}
\usepackage[hidelinks]{hyperref}

\title{A learnable Haar tokenizer for ImageNet-100\\
\large Architecture, training, and ablation results from \texttt{haar\_tok.py}}
\author{}
\date{}

\begin{document}
\maketitle

\begin{abstract}
\texttt{haar\_tok.py} trains a classifier whose first stage is a four-level
Haar analysis of a $256\times256$ image. With the filters held at their Haar
initialisation, each $16\times16$ image patch is replaced by the $768$
coefficients of an orthonormal transform of that patch. Those coefficients are
then mixed by a twelve-block network and pooled into one of $100$ ImageNet-100
classes. The run \texttt{tok\_learn\_conv\_100} trains the filters, keeps a
depthwise-convolution mixer, and runs for $100$ epochs. Its best validation
checkpoint reaches \textbf{73.80\% top-1} on the official test split of
\textbf{5{,}000 images} (validation 78.09\% on 13{,}000 images). Thirty-epoch
ablations show that a convolution mixer beats attention, that unfreezing the
Haar filters changes accuracy only slightly at $30$ epochs, and that appending
detail magnitudes is the strongest $30$-epoch variant.
\end{abstract}

\section{What the script writes}

A run named \texttt{--name} (default \texttt{haar\_tok}) creates the directory
\texttt{training/checkpoints/runs/imagenet100/<name>/} and writes:

\begin{center}
\begin{tabular}{@{}ll@{}}
\toprule
File & Contents \\
\midrule
\texttt{args.json} & every command-line setting for the run \\
\texttt{history.csv} & one row per epoch: losses, top-1, learning rate, orthogonality \\
\texttt{last.pt} & weights and AdamW state after the latest epoch \\
\texttt{best.pt} & the same state at the best validation top-1 so far \\
\texttt{summary.json} & parameter count, best validation top-1, and the single test evaluation \\
\bottomrule
\end{tabular}
\end{center}

The test number below is not the accuracy of epoch $100$. After training, the
script reloads \texttt{best.pt} and evaluates that checkpoint once. For
\texttt{tok\_learn\_conv\_100} the best validation epoch is epoch $95$.

\section{Data}

The dataset is the Hugging Face release \texttt{ilee0022/ImageNet100}, read
from the cache under \texttt{training/checkpoints/data/imagenet100}. Its
recorded splits are:

\begin{center}
\begin{tabular}{@{}lr@{}}
\toprule
Split & Images \\
\midrule
Train & 117{,}000 \\
Validation & 13{,}000 \\
Test & 5{,}000 \\
\bottomrule
\end{tabular}
\end{center}

There are $100$ classes. The training loader uses batch size $256$ and drops
the last incomplete batch, so each epoch updates on $116{,}992$ of the
$117{,}000$ training images. Validation and test loaders do not drop images.
The test loader uses batch size $512$. Thus $73.80\%$ test accuracy is
$3{,}690$ correct predictions out of $5{,}000$.

Training images are randomly resized and cropped to $256\times256$ with scale
range $(0.67,1)$, flipped horizontally, passed through RandAugment with two
operations of magnitude $9$, normalised with the ImageNet channel statistics
$(0.485,0.456,0.406)$ and $(0.229,0.224,0.225)$, and then randomly erased with
probability $0.25$. Evaluation resizes the short side to $293$
($256/0.875$, rounded) and takes a $256\times256$ centre crop before the same
normalisation. No augmentation is used at evaluation.

\section{The Haar tokenizer}

\subsection{One analysis step}

The tokenizer applies a stride-$2$ convolution with four $2\times2$ kernels
and no padding. Because the support is exactly one $2\times2$ block, neighbouring
blocks do not overlap. The four kernels are initialised to the orthonormal
two-dimensional Haar bank
\begin{equation}
\begin{aligned}
W_{\mathrm{LL}} &= \tfrac12\begin{bmatrix} 1 & 1 \\ 1 & 1 \end{bmatrix}, &
W_{\mathrm{LH}} &= \tfrac12\begin{bmatrix} 1 & -1 \\ 1 & -1 \end{bmatrix}, \\[4pt]
W_{\mathrm{HL}} &= \tfrac12\begin{bmatrix} 1 & 1 \\ -1 & -1 \end{bmatrix}, &
W_{\mathrm{HH}} &= \tfrac12\begin{bmatrix} 1 & -1 \\ -1 & 1 \end{bmatrix}.
\end{aligned}
\end{equation}
Flattening each kernel into a row gives a $4\times4$ matrix $W$. These rows
are orthonormal: $WW^{\mathsf T}=I_4$. Consequently one analysis step replaces
every $2\times2$ neighbourhood by four coefficients and preserves Euclidean
norm. The same four kernels are used independently on red, green, and blue.
There is one bank per wavelet level, not one bank per channel.

Low-pass coefficients are passed to the next level. The three detail bands at
that level are retained.

\subsection{Four levels on a $16\times16$ patch grid}

The image is $256\times256$ and the patch size is $16$, so there are four
dyadic levels and a $16\times16$ grid of $256$ tokens. After level $\ell$
(starting at $\ell=0$) the detail maps have spatial size $256/2^{\ell+1}$.
Those maps are rearranged with pixel unshuffle so that every coefficient whose
receptive field lies inside a given $16\times16$ image region is stored at that
region's token. For one colour channel the coefficient counts are:

\begin{center}
\begin{tabular}{@{}lrrr@{}}
\toprule
Level & Spatial size before packing & Detail coefficients & Role \\
\midrule
0 & $128\times128$ & $3\times8\times8=192$ & finest edges and texture \\
1 & $64\times64$ & $3\times4\times4=48$ & \\
2 & $32\times32$ & $3\times2\times2=12$ & \\
3 & $16\times16$ & $3\times1\times1=3$ & coarsest details \\
3 (LL) & $16\times16$ & $1$ & local average of the patch \\
\midrule
\multicolumn{2}{@{}l}{Total per channel} & \multicolumn{2}{r@{}}{$256=16^{2}$} \\
\bottomrule
\end{tabular}
\end{center}

Three channels therefore produce
\begin{equation}
3\times16\times16=768
\end{equation}
numbers at every token. While the filters remain the Haar bank above, this
vector is an orthonormal reparameterisation of the $768$ pixels in the
$16\times16\times3$ patch: nothing inside the patch is discarded, and no pixel
from a neighbouring patch enters the token. This is the property a length-$2$
Haar filter has and a longer wavelet filter does not.

\subsection{Making the filters trainable}

In the $74\%$ run the $64$ filter values ($4$ levels $\times4$ kernels
$\times2\times2$) are parameters. They are computed in float32 even though the
rest of the network runs in bfloat16. A soft penalty discourages them from
leaving the orthonormal family. If $W_\ell\in\mathbb R^{4\times4}$ is the
flattened bank at level $\ell$,
\begin{equation}
\mathcal L_{\mathrm{ortho}}
= \frac1{4}\sum_{\ell=0}^{3}\lVert W_\ell W_\ell^{\mathsf T}-I_4\rVert_F^2.
\end{equation}
The training loss adds $10^{-2}\,\mathcal L_{\mathrm{ortho}}$ to cross-entropy.
The recorded penalty itself is about $3\times10^{-3}$ after epoch $1$ and about
$0.214$ at epoch $95$, so the learned filters do drift. The weighted penalty
contributes roughly $0.002$ to a cross-entropy near $0.90$ late in training, which
is too small to freeze them at Haar. Filter weights are also excluded from
AdamW weight decay, together with the positional embedding and every
one-dimensional parameter.

The fixed-filter ablation registers the same initial bank as a buffer. Its
penalty is identically zero and it has $64$ fewer trainable parameters
($6{,}768{,}996$ instead of $6{,}769{,}060$).

An optional modulus mode appends the absolute value of every detail map before
the linear embedding. The low-pass map is not duplicated. The resulting token
width is $1{,}533$ and the recorded model has $6{,}966{,}430$ parameters. The
$100$-epoch run does not use this mode.

\section{From tokens to a class}

Each $768$-dimensional token is layer-normalised and linearly mapped to
dimension $d=256$. A learned positional vector is added. There is no class
token. The sequence of $256$ tokens passes through $12$ pre-norm residual
blocks. Block $i$ uses drop-path rate $0.1\cdot i/11$ and two layer-scale
vectors initialised to $10^{-4}$:
\begin{equation}
\begin{aligned}
z &\leftarrow z + \gamma_1\odot \mathrm{Mixer}(\mathrm{LN}(z)),\\
z &\leftarrow z + \gamma_2\odot \mathrm{MLP}(\mathrm{LN}(z)).
\end{aligned}
\end{equation}
The MLP expands $256$ to $1{,}024$, applies GELU, and returns to $256$.

The mixer used for the $74\%$ model is a depthwise $7\times7$ convolution. The
sequence is reshaped to a $16\times16$ feature map, convolved with one
$7\times7$ kernel per channel, and reshaped back. It mixes only neighbouring
tokens. The attention ablation instead uses four-head scaled dot-product
attention over all $256$ tokens, with a combined query-key-value projection.
That model has $9{,}773{,}476$ parameters.

After the last block the tokens are layer-normalised and averaged over the
$256$ positions. A linear layer maps the $256$-dimensional mean to $100$
logits. Softmax is not applied inside the model.

\section{Optimisation}

The $100$-epoch configuration recorded in
\texttt{tok\_learn\_conv\_100/args.json} is:

\begin{center}
\begin{tabular}{@{}ll@{}}
\toprule
Choice & Value \\
\midrule
Epochs & $100$ \\
Warmup & $5$ epochs \\
Batch size & $256$ \\
Optimiser & fused AdamW \\
Peak learning rate & $10^{-3}$ \\
Schedule & linear warmup, then cosine decay to zero, updated every step \\
Weight decay & $0.05$ on matrix weights only \\
Label smoothing & $0.1$ \\
Orthogonality weight & $10^{-2}$ \\
Precision & bfloat16 autocast, tokenizer in float32 \\
Seed & $42$ \\
\bottomrule
\end{tabular}
\end{center}

The thirty-epoch ablations use the same recipe except that warmup is $3$
epochs rather than $5$. Train top-1 is measured under the smoothed training
loss; reported train loss is the smoothed cross-entropy and does not include
the orthogonality term. Validation and test losses are unsmoothed
cross-entropy. Checkpoint selection uses validation top-1 only. The test split
is not read until training has finished.

\section{The $100$-epoch result}

\texttt{tok\_learn\_conv\_100} is learnable Haar filters, the depthwise
convolution mixer, no modulus features, and $100$ epochs.

\begin{center}
\begin{tabular}{@{}lrr@{}}
\toprule
 & Epoch 95 (selected) & Epoch 100 (last) \\
\midrule
Train loss & 0.897 & 0.895 \\
Train top-1 & 96.70\% & 96.79\% \\
Validation loss & 0.998 & 0.997 \\
Validation top-1 & 78.09\% & 78.01\% \\
\bottomrule
\end{tabular}
\end{center}

Epoch $95$ is the best of the $100$ validation measurements: $10{,}152/13{,}000$.
Reloading that checkpoint and evaluating the untouched test split gives

\begin{equation}
\text{test top-1}=73.80\%,\qquad \text{test loss}=1.237.
\end{equation}

The train accuracy near $97\%$ against $78\%$ validation and $74\%$ test shows
that the model has memorised much of the training set by the time the cosine
schedule has decayed. The orthogonality penalty at epoch $95$ is $0.214$.

\section{Ablations}

All ablations below were selected the same way: best validation checkpoint,
then one test evaluation. The first four rows use $30$ epochs. The last row is
the run above.

\begin{center}
\small
\begin{tabular}{@{}lccccrrrr@{}}
\toprule
Run & Epochs & Filters & Mixer & $|{\cdot}|$ & Params & Best val & Test \\
\midrule
\texttt{tok\_fixed\_conv} & 30 & fixed Haar & conv $7\times7$ & no & 6.769M & 73.73\% & 69.54\% \\
\texttt{tok\_learn\_conv} & 30 & learned & conv $7\times7$ & no & 6.769M & 74.23\% & 69.46\% \\
\texttt{tok\_learn\_attn} & 30 & learned & attention & no & 9.773M & 70.24\% & 65.62\% \\
\texttt{tok\_learn\_mod} & 30 & learned & conv $7\times7$ & yes & 6.967M & 75.92\% & 72.52\% \\
\texttt{tok\_learn\_conv\_100} & 100 & learned & conv $7\times7$ & no & 6.769M & 78.09\% & \textbf{73.80\%} \\
\bottomrule
\end{tabular}
\end{center}

At equal epoch count, replacing the convolution mixer with global attention
costs about $4$ test points and adds about three million parameters. Unfreezing
the Haar filters improves validation by half a point at $30$ epochs and does
not improve test accuracy ($69.46\%$ learned versus $69.54\%$ fixed). The gain
that produces $73.80\%$ is therefore mostly the additional $70$ epochs, not the
$64$ extra filter parameters. Adding detail magnitudes is the strongest change
available inside the $30$-epoch budget: $72.52\%$ test, within $1.3$ points of
the $100$-epoch model, without a matching $100$-epoch modulus run.

These tokenizer runs are a different model from the earlier operator-mixer
experiments \texttt{haar\_orth} and \texttt{haar\_over}. Those runs reconstruct
full-resolution Haar subspaces and mix channels with $D\times D$ matrices.
After $30$ epochs their best validation accuracies were $27.71\%$ and
$30.27\%$. They are not the source of the $73.80\%$ test result.

\section{End-to-end shape}

For one image in the $74\%$ configuration:
\begin{equation}
\begin{aligned}
\mathbb R^{3\times256\times256}
&\;\xrightarrow{\;\text{4-level Haar}\;}\;
\mathbb R^{768\times16\times16}\\
&\;\xrightarrow{\;\text{LN, linear, position}\;}\;
\mathbb R^{256\times256}\\
&\;\xrightarrow{\;\text{12 mixer blocks}\;}\;
\mathbb R^{256\times256}\\
&\;\xrightarrow{\;\text{mean over tokens}\;}\;
\mathbb R^{256}
\;\xrightarrow{\;\text{linear}\;}\;
\mathbb R^{100}.
\end{aligned}
\end{equation}
The middle $256$ is the token count, not the channel count. Each of those
tokens is the multi-scale Haar description of one non-overlapping $16\times16$
patch, projected from $768$ coefficients to $256$ features.

\end{document}
